%% ma: L11-B16-0011
%% lop: 11
%% chuong: 5
%% bai: 16
%% dang: 11.16.6
%% loai: TL
%% muc: TH
%% dap_an: "-"
%% nguon: "tai-lieu-hoc-tap-mon-toan-lop-11-theo-chuong-trinh-giao-duc-pho-thong-2018.pdf (Lớp Toán Thầy Duy, Phạm Lê Duy, THPT Chu Văn An), Chương 3 Bài 2, Dạng 3, Ví dụ 11 (tr.357)"
%% kien_thuc: "Lớp 11 Bài 16 (giới hạn của hàm số)"
%% trang_thai: cho-duyet
%% ly_do: "Dạng mới đề xuất (11.16.x), chờ giáo viên duyệt dạng; tài liệu gốc không có lời giải mẫu, đáp án do kho tự giải và kiểm bằng sympy."
\begin{ex}
Tính giới hạn của các hàm số sau.

a) $f(x)=\begin{cases}\dfrac{2x}{\sqrt{1-x}}&\text{với }x<1\\\sqrt{3x^2+1}&\text{với }x\ge1\end{cases}$. Tính $\lim\limits_{x\to1^+}f(x)$.

b) $f(x)=\begin{cases}\dfrac{x^2+1}{1-x}&\text{với }x<1\\\sqrt{2x-2}&\text{với }x\ge1\end{cases}$. Tính $\lim\limits_{x\to1^-}f(x)$.
\loigiai{\sloppy a) Khi $x\to1^+$ thì $x>1$ nên $f(x)=\sqrt{3x^2+1}\to\sqrt4=2$.

b) Khi $x\to1^-$ thì $x<1$, $f(x)=\dfrac{x^2+1}{1-x}$ có tử $\to2>0$, mẫu $\to0^+$, nên giới hạn bằng $+\infty$.}
\end{ex}
